En este trabajo se presentarán los resultados obtenidos de estudiar el comportamiento experimental de distintos algoritmos de ordenamiento y multiplicación de matrices. Los algoritmos de ordenamiento que serán estudiados en este trabajos son: merge sort, quick sort, patience sort y el algoritmo sort de la librería estandar de C++.Por otro lado, también, los algoritmos de multiplicación de matrices que estudiaremos son: naive y strassen. \\
El objetivo es medir el tiempo de ejecución y la memoria utilizada por cada algoritmo bajo distintos tamaños de entrada, tipos de datos y dominios de valores, para contrastar el comportamiento observado empíricamente con lo que predice su complejidad teórica